\section{Approximation hardness of \textsc{Connectivity}}
In this section, we show that \textsc{Connectivity} cannot be approximated within a factor of $O\br{\log^{2-\epsilon} n}$ for any $\epsilon > 0$, unless $\text{NP} \subseteq \text{ZTIME}(n^{\text{polylog}(n)})$. To achieve this result we give a reduction from the \textsc{Group Steiner Tree} (\textsc{GST}) problem, which is known to be hard to approximate within a factor of $O\br{\log^{2-\epsilon} n}$~\cite{PolylogarithmicInapproximability}. In \textsc{GST}, we are given a graph $G = (V, E)$, a cost function $c: E \to \mathbb{R}_{\geq 0}$, and a collection of groups $g_1, \dots, g_k \subseteq V$. The goal is to find a subtree $T$ of $G$ that contains at least one vertex from each group minimizing $c(T)=\sum_{e\in E\br{T}}c\br{e}$. 

In order to prove this claim we need to use the specific properties of the construction used in order to show the hardness of approximation of \textsc{GST} given in~\cite{PolylogarithmicInapproximability}. 
Let $m$ be the size of some PCP system for $\text{NP}$. The instance of \textsc{GST} obtained in the reduction from this PCP system is a tree with the following properties:
\begin{enumerate}
    \item $k=m^{O\br{\log m}}$.
    \item The tree is a complete $d$-ary tree of height $H=O\br{\log^{\beta} m}=O\br{\log^{\frac{\beta}{\beta+2}} k}$ for some $\beta > 0$.
    \item The number of vertices of the tree is $O\br{d^H} = m^{O\br{\log^c m}}$ for some constant $c$.
    \item The cost of each edge is exactly $1/\tau$ times the cost of its parent edge for any desired constant $\tau > 1$.
\end{enumerate}
We call a tree of above type \emph{elegant}. We have:
\begin{theorem}[\cite{PolylogarithmicInapproximability}]
    \textsc{GST} cannot be approximated within a factor of $O\br{\log^{2-\epsilon} n}$ for any $\epsilon > 0$, unless $\text{NP} \subseteq \text{ZTIME}(n^{\text{polylog}(n)})$, even if the input graph is an elegant tree.
\end{theorem}
\begin{theorem}
    \textsc{Connectivity} cannot be approximated within a factor of $O\br{\log^{2-\epsilon} n}$ for any $\epsilon > 0$, unless $\text{NP} \subseteq \text{ZTIME}(n^{\text{polylog}(n)})$, even if the input graph is a DAG.
\end{theorem}
\begin{proof}
    The idea of our construction is as follows: The crucial difficuilty of the reduction is the fact that the cost criterion in \textsc{GST} is the sum of the costs of the edges, while in \textsc{Connectivity} it is the maximum cost of labels activated per vertex. In our construction, a desirable goal is for each edge outgoing of a vertex $v$ to represent a cost of all of the edges in the subtree rooted at the corresponding child of $v$. To achieve this, we will introduce a significant number of child representatives, each encoding a particular choice of cost of edges in the subtree rooted at $v$. In fact, each vertex will encode even richer information, i. e., the number of edges chosen at any level of the subtree rooted at $v$. Since the input tree is elegant, this information exatly determines the cost of the subtree rooted at $v$. Then, the next challenge is to force any reasonable algorithm to never exceed the cost entrusted to a given subtree by a too large factor. To achieve this, we will scale the costs in all subtrees using an appopriate (and usually differing) factor, so that each vertex in the tree has the same "allowed" cost. At last, we will also need to guess the value of the optimal solution to be able to scale the costs in the right way. Since, the number of possible values of the optimal solution is quasipolynomial this will not affect the running time.

    Assume, that we have guessed the value of the optimal solution $B$. Let $T$ be the input tree of \textsc{GST}. We construct a DAG, $D_B$ as follows: In the first stage of the construction  we assume $D_B$ is a rooted tree. Let $r$ be the root of $T$. Let $b_r=B$. We set $r'$ as the root of $D_B$ and we entrust it with a budget of $b_r$. The rest of the construction is recursive. The recurrence consists of a leaf $\ell$ of $D_B$ and a (some) root vertex $r$ of a subtree of $T$. For any child of $r$, call it $v$, we create $t_v=\spr{V\br{T_v}}^{O\br{H\br{T_v}}}$ vertex representatives $u_i$ of $v$, for $i\in [t_v]$ if the following conditions hols: For any such representative let $\brc{\mathcal{N}_{u_i}\br{j}}_{j=1}^{H\br{T_v}}$ be a sequence of non-negative integers such that $\mathcal{N}_{u_i}\br{j}\in \brc{0,\ldots, \spr{V\br{T_{v}}}}$ and there exists at least one subtree of $T_v$ with that number of vertices at each level. Note that this can be checked in quasipolynomial time, since the tree is elegant. Let $b_{u_i}=\sum_{j=1}^{H\br{T_v}}\mathcal{N}_{u_i}\br{j}\cdot \tau^j$ denote the budget entrusted to $u_i$. We demand that $b_{u_i}\leq b_{\ell}$. 
    
    Observe, that if an algorithm were to pick exactly the number of edges at each level of the subtree rooted at $v$ as encoded by $\mathcal{N}_{u_i}$, then the cost of the subtree rooted at $v$ would be exactly $b_{u_i}$. We connect the leaf $\ell$ to $u_i$ (with a directed edge) and create a label $u_i$ which we add to $\lambda\br{\ell}$ and $\lambda\br{u_i}$. We also set the cost of $u_i$ to be $b_{u_i}\cdot B/b_{\ell}$ for $\ell$ and $0$ for $u_i$. In this way we ensure that if the subtree $T_{\ell}$ is entrusted with a budget of $b_{\ell}$, then exceeding this budget $\alpha$ times in the original instance would correspong to paying $\alpha\cdot B$ in the new instance.
    
    In the second part of the construction, we need to encode the groups $g_1, \dots, g_k$. We assume that there is only one pair of source and sink subset in the \textsc{Connectivity} instance. The source conists solely of $r'$. For each group $g_i$ we create a vertex $s_i$ and we connect all the representatives of the vertices in $g_i$ to $s_i$ (directed). We also add terminal $g_i$ to both $\lambda\br{s_i}$ and all the representatives of the vertices in $g_i$. We set the cost of $g_i$ to be universally $0$. Finally, we set the sink subset to be $\brc{s_1, \dots, s_k}$.
    \begin{lemma}
        The size of $D_B$ is quasipolynomial in $m$.
    \end{lemma}
    \begin{proof}
        First, the number of sink nodes is $k = m^{O(\log m)}$. It remains to show that the number of representative vertices is quasipolynomial. Firstly, we show that each vertex has a quasipolynomial number of children. Let $r$ be such a vertex. Then, for each child vertex $v_i$, for $i\in[d]$, $r$ has $t_{v_i}=\spr{V\br{T_{v_i}}}^{O\br{H\br{T_{v_i}}}}$ children. Therefore the overall number of children is bounded by:
        \begin{align*}
            \sum_{i=1}^{d} t_{v_i} = \sum_{i=1}^{d} \spr{V\br{T_{v_i}}}^{O\br{H\br{T_{v_i}}}} \leq \sum_{i=1}^{d} \spr{V\br{T}}^{O\br{H\br{T}}} = d\cdot \spr{V\br{T}}^{O\br{H\br{T}}} = m^{O\br{\log^c m}}
        \end{align*}
        Since the depth of the DAG is at most the depth of the tree (plus $1$), we know that there exists an apropriate constant $c'$ such that the number of vertices in $D_B$ is at most \begin{align*}
            \sum_{i=0}^{H\br{T}+1} \br{\spr{V\br{T}}^{O\br{H\br{T}}}}^i = \sum_{i=0}^{H\br{T}+1} m^{O\br{\log^c m}\cdot i}  = \frac{m^{O\br{\log^c m}\cdot \br{H\br{T}+1}} - 1}{m^{O\br{\log^c m}} - 1} = m^{O\br{\log^{c'} m}}
        \end{align*}
        where the second inequality is by the fact that the sum is a geometric series and the last equality is by the fact that $H\br{T}=O\br{\log^{\beta}{m}}$.
    \end{proof}
    Now, we wish to show that upon guessing the value $B$, such that the optimal solution of the \textsc{GST} instance is $B$, any $\alpha$-approximation algorithm for \textsc{Connectivity} on $D_B$ would give an $\alpha$-approximation for the \textsc{GST} instance. 
        
    Firstly, without loss of generality we may assume that the algorithm picks at most one representative for each vertex of $T$. Indeed, if it picks more than one representative for a vertex $v$, then we can replace all the representatives of $v$ with a single representative $u$. Indeed, let $u_1,\dots, u_p$ be the representatives of $v$ picked by the algorithm. Let $\mathcal{N}_u\br{j}=\min\brc{d^j, \sum_{i=1}^{p} \mathcal{N}_{u_i}\br{j}}$ for each level $j$. It is easy to see that: Firstly, the cost of $u$ is at most the cost of all the representatives of $v$ picked by the algorithm. Secondly, the $\mathcal{N}_u\br{j}$ suffices to pick all edges in subtree rooted at $v$ that are picked by the representatives $u_1,\dots, u_p$. Therefore, we can replace all the representatives of $v$ with $u$ without increasing the cost of the solution and without losing any edge in the subtree rooted at $v$ that is picked by the algorithm.
        
    Now, we show how to recover the solution. We begin by removing all of the sink vertices from $D_B$ to obtain a tree. Initialize $\mathcal{T}$ to be an empty tree which will be the solution for \textsc{GST}. Then we traverse the tree in a top-down manner, starting from the root. Let $r$ be some vertex met along the way. Whenever an edge to a representative $u$ of a child of $r$, $v$ is picked, we append $uv$ to the solution of \textsc{GST} and we continue the traversal towards $u$. It is easy to see that the solution created in this way indeed contains at least one vertex from each group and thus is a feasible solution for \textsc{GST}. 
    \begin{lemma}
        If the optimal solution of the \textsc{GST} instance is $B$, then if there exists an $\alpha$-approximation algorithm for \textsc{Connectivity} on $D_B$, then the cost of $\mathcal{T}$ is at most $\alpha\cdot B$.
    \end{lemma}
    \begin{proof}
        We show this by a bottom-up induction over the tree $D_B$ (without the sink vertices). We claim, that 
        
        It is easy to see that the claim holds for the leaves of $D_B$. This is because each leaf $\ell$ of $D_B$ is a singular existing representative of some leaf of $T$ and thus, each edge of its parent 

    \end{proof}
\end{proof}
